\scanpage{121}
\mspage{2.1}
2. Ramification Index

Before proceeding to the proofs of the main results on the ramification
index of a local $R$-algebra $S$ in the case that $R$ and $S$ are \inserted{integrally closed} domains,
we establish some of the more formal properties of $r(S/R)$.

Lemma 2.1: Let $S$ be a local $R$-algebra and $T$ a local $S$-algebra.
If $S$ is an unramified $R$-algebra, then the homomorphism $T_R^f\longrightarrow T_S^f$ induced
from the natural map $T\otimes_R T\longrightarrow T\otimes_S T$ is an isomorphism. In particular
if $T$ is unramified over $R$, then $T_R^f\approx T_T^f=T$.

\uline{Proof}: Since $S$ is an unramified $R$-algebra, the exact sequence
$0\longrightarrow\unsure{\mathfrak g|\mathfrak J}\longrightarrow S\otimes_R S\longrightarrow S\longrightarrow0$ splits as $S^e$-modules ([5]). Therefore
$0\longrightarrow T_R^e\otimes_{S^e}\unsure{\mathfrak g|\mathfrak J}\longrightarrow T_R^e\longrightarrow T_R^e\otimes_{S^e}S\longrightarrow0$ is exact and splits as $T_R^e$-
modules. By one of the standard associativity laws we have that $T_R^e\otimes_{S^e}S=
(T\otimes_R T)\otimes_{S\otimes S}S=T\otimes_S T=T_S^e$. Thus we find that the natural epimorphism
$T_R^e\longrightarrow T_S^e$ splits as $T_R^e$-modules. Let $\underline N$ be the maximal ideal of $T$ and $\unsure{\eta|\mathfrak N}$ the
preimage of $N$ under the map $T\otimes_R T\longrightarrow T$. Localizing by $\unsure{\eta|\mathfrak N}$, we see that
$T_R^f\longrightarrow T_S^f$ splits as $T_R^f$-module since $(T_R^e)_{\unsure{\eta|\mathfrak N}}=T_R^f$ and $(T_S^e)_{\unsure{\eta|\mathfrak N}}=T_S^f$ also.
Consequently $T_R^f=T_S^f$ since $T_R^f$ is a local ring.

Lemma 2.2: Let $S$ be a local $R$-algebra. Then we have $\widehat{S^f}\approx\widehat S^f$ where $\wedge$
denotes the completion of a local ring under the topology given by the
powers of the maximal ideal.

\uline{Proof}: Since $S$ is a local $R$-algebra, $\widehat S$ is a finite $\widehat R$-module and hence
$\widehat S^e=\widehat S\otimes_{\widehat R}\widehat S$ is a complete semi-local ring under $\underline m\widehat S^e$-adic topology were $\underline m$ is
the maximal ideal in $R$. Therefore $\widehat S_{\widehat R}^f$, the localization of $\widehat S^e$ by a maximal
ideal, is a complete local ring and hence the canonical homomorphism
$S^f\longrightarrow\widehat S^f$ induces the homomorphism $\widehat{S^f}\longrightarrow\widehat S^f$. Since $\widehat{S^f},\widehat S^f$ are both
\mspage{120}
